\documentclass[12pt,a4paper]{article}

\usepackage[utf8]{inputenc}
\usepackage[T1]{fontenc}
\usepackage[spanish]{babel}
\usepackage{lmodern}
\usepackage{microtype}
\usepackage{geometry}
\usepackage{setspace}
\usepackage{parskip}

\geometry{
  top=2.5cm,
  bottom=2.5cm,
  left=3cm,
  right=3cm
}

\setstretch{1.15}

\title{¿Dónde localizar la relación principal entre neoliberalismo e IA?}
\author{Vicente Domínguez Arca}
\date{}

\begin{document}

\maketitle

Con esta nueva entrada trato de buscar un contacto más concreto con el posible lector. La pregunta que la encabeza es provocadora y alguien puede pensar que este texto va a dirigirse hacia una crítica o exposición de algo tan evidente como la inserción de la inteligencia artificial en la sociedad mediante empresas privadas y, muy particularmente, mediante grandes empresas privadas. No. No me voy a dirigir ahí. Tampoco intentaré desarrollar aquí un manifiesto sobre una inteligencia artificial socializada. Me interesa algo mucho más fundamental, algo que considero necesario plantear en estos tiempos y, especialmente, en esta actualidad en la que proliferan muchos, y a veces demasiados, comentarios sobre futuribles implicados en y con la inteligencia artificial.

Aunque esta entrada estará marcada y atravesada por una fuerte cuestión epistemológica y por una clara tendencia filosófica, comenzaré desde un lugar mucho más concreto, como investigador en activo que programa y trabaja con aplicaciones de redes neuronales. Sin entrar en detalles técnicos innecesarios, quizá al lector le suene, o incluso le resulte tremendamente familiar por sus propios conocimientos sobre el tema, una de las partes cruciales en el desarrollo de cualquier red neuronal basada en aprendizaje a partir de datos: la etapa de entrenamiento. Básicamente, partiendo de una configuración determinada de algoritmos matemáticos, de los que no hablaremos aquí, necesitamos ajustar determinados parámetros utilizando datos de entrenamiento. Si deseamos que una red neuronal proporcione texto generado, debemos entrenarla con grandes cantidades de texto. Si buscamos que sea capaz de detectar masas tumorales en imágenes radiológicas, deberemos entrenarla con imágenes adecuadas para esa tarea, junto con la información necesaria para que el sistema pueda ajustar sus parámetros a aquello que buscamos reconocer. Y así sucesivamente.

Aquí quiero detenerme para que el lector interiorice algo que, quizá precisamente por parecer evidente, demasiado a menudo se pasa por alto, incluso, y más crucialmente, en los ámbitos investigador y técnico. Cualquiera que sea el origen de los datos con los que se entrene una red neuronal, esos datos no pertenecen al Universo en un sentido absoluto, pertenecen al acceso humano al Universo. Podremos discutir cuánto han sido preprocesados, podremos argumentar largo y tendido que determinados datos son verdaderamente crudos, podremos automatizar su adquisición hasta reducir extraordinariamente la intervención humana inmediata, pero seguirán siendo datos de nuestro \emph{Huniverso}.

Llamo \emph{Huniverso} a aquella parte del Universo a la que el ser humano tiene acceso, bien mediante sus sentidos naturales, bien mediante su capacidad de procesamiento y relación, bien mediante las extensiones tecnológicas que ha construido para ampliar sus posibilidades de observación. El Huniverso no es el Universo, sino el Universo accesible desde nuestra condición humana. No disponemos de datos del Universo en sí, disponemos de datos del Huniverso. Y eso significa que todo dato accesible para nosotros contiene inevitablemente una mediación humana, no necesariamente porque una persona haya manipulado deliberadamente cada registro, sino porque su propia condición de dato disponible para nosotros implica ya una selección, una parametrización, una medición, una codificación o, cuando menos, una forma de acceso determinada por aquello que somos y por aquello que podemos llegar a percibir o registrar.

Podemos decir, por ejemplo, que una red neuronal se entrena con datos de ratones, pero conviene ser precisos. No son datos producidos desde el acceso del ratón al Universo, ni representaciones elaboradas desde aquello que un ratón percibe, distingue o considera relevante. Son datos humanos acerca de aquello que nosotros damos en llamar ratones: imágenes obtenidas por dispositivos diseñados por humanos, variables escogidas por humanos, señales registradas mediante instrumentos construidos por humanos, comportamientos clasificados desde categorías humanas y parámetros definidos dentro de nuestros sistemas de observación. Podemos llamar a eso datos de ratones por comodidad, pero la expresión no debe hacernos olvidar que continúan siendo material huniversal.

Esta apreciación resulta importante, mucho, tal vez sea crucial, porque cualesquiera que sean las redes neuronales que configuremos mediante algoritmos que deban ajustarse a datos dispuestos por seres humanos, sus resultados quedarán necesariamente relacionados con esa realidad huniversal desde la que aquellos datos fueron obtenidos. Existe, naturalmente, una réplica inmediata: una red neuronal puede encontrar correlaciones en sus datos de entrenamiento que ningún ser humano haya percibido. Sin duda. Y puede hacerlo extraordinariamente rápido. Cuanto más se definan y orienten los algoritmos hacia una tarea determinada, y cuanto más extensos y adecuados sean los datos de entrenamiento, su capacidad para encontrar estructuras, regularidades y correlaciones puede resultar inigualable para un individuo. Puede procesar volúmenes de información inimaginables para un ser humano, puede encontrar relaciones que nosotros no habíamos advertido y puede producir resultados que nos sorprendan. Pero incluso entonces esas correlaciones procederán del Huniverso, porque aquello sobre lo que correlaciona ha llegado hasta el entrenamiento codificado mediante datos accesibles a seres humanos o mediante mecanismos informáticos conocidos, diseñados y construidos desde nuestra capacidad tecnológica.

En este sentido, las redes neuronales actuales pueden ser más rápidas que un ser humano, pueden ser más eficaces en la búsqueda de determinadas correlaciones, pueden tratar volúmenes de datos inaccesibles para cualquier individuo y pueden encontrar regularidades que no habíamos observado. Pero no pueden, por construcción, salir del espacio informacional al que han tenido acceso. No pueden recibir aquello que nunca ha sido percibido, medido, registrado, parametrizado o representado desde nuestro Huniverso. Pueden recombinar, correlacionar, generalizar y producir resultados nuevos dentro de ese espacio, pero no pueden acceder a un afuera que jamás haya sido incorporado de ningún modo a los datos o a los mecanismos con los que han sido construidas.

Llegados aquí aparece una de las cuestiones que más preocupación, sentida o artificiosa, discusiones aparte, ha provocado alrededor de la inteligencia artificial: la posibilidad de una futura extinción humana debida a la acción de inteligencias artificiales suficientemente complejas y superiores a nosotros en determinadas capacidades. Bien, planteémonos este problema desde lo que acabamos de exponer. El ser humano muestra una persistente tendencia hacia la supervivencia, procura sobrevivir, reproducirse y conservar las condiciones que hacen posible su continuidad. Naturalmente existen excepciones individuales, conductas autodestructivas, conflictos, guerras y episodios históricos en los que los seres humanos han puesto en peligro a otros seres humanos e incluso las condiciones de su propia continuidad, pero, considerada la especie en su conjunto, resulta difícil negar la presencia recurrente de esa orientación hacia la supervivencia. Es, además, algo que observamos en nuestro Huniverso en aquello que damos en llamar otras especies.

Si esa tendencia atraviesa profundamente nuestra existencia, tiene necesariamente que haber dejado su huella en aquello que producimos. En nuestra literatura, en nuestro arte, en nuestras culturas, en nuestras instituciones, en nuestra ciencia, en nuestras tecnologías, en nuestros sistemas simbólicos y en nuestra manera de modelar aquello que sentimos y observamos. Y, por extensión, en los datos con los que entrenamos nuestras redes. De ahí que este autor encuentre dificultades para justificar que una red neuronal entrenada exclusivamente mediante material huniversal pueda desarrollar, por sí misma y como finalidad propia, una orientación hacia la desaparición de la especie humana. No porque una inteligencia artificial no pueda participar causalmente en una catástrofe, no porque un sistema automatizado no pueda adoptar decisiones extremadamente perjudiciales, ni mucho menos porque la inteligencia artificial carezca de peligros. La cuestión es otra: ¿de dónde procedería esa supuesta finalidad autónoma de exterminio si todo aquello a lo que el sistema ha tenido acceso está atravesado, en mayor o menor medida, por una realidad huniversal en la que la supervivencia aparece de forma constante como una de las grandes tendencias de los seres vivos que observamos?

No pretendo convertir esa pregunta en una demostración matemática ni cerrar con ella el debate sobre los riesgos de la inteligencia artificial. Muy al contrario. Según este autor, la inteligencia artificial alberga enormes peligros y retos, muchos de los cuales probablemente todavía están por llegar. Pero quizá algunos de esos peligros no se encuentren donde más cómodamente imaginamos. Porque los datos humanos no contienen solamente nuestra tendencia hacia la supervivencia. Contienen también nuestros otros sesgos, nuestras estructuras sociales, nuestras desigualdades, nuestras instituciones, nuestras guerras, nuestras formas de poder y, naturalmente, nuestras formas de acumulación.

Aquí comienza a aparecer la pregunta que da título a esta entrada. Una enorme proporción de la información digital disponible actualmente ha sido producida, recopilada, seleccionada, almacenada y procesada durante las últimas décadas. Naturalmente, gran parte del material al que esos registros hacen referencia puede proceder de tiempos mucho más antiguos, incluso de milenios, pero su transformación en datos digitalmente accesibles, clasificables, procesables y entrenables ha ocurrido de manera muy particular dentro de sociedades tecnológicamente desarrolladas y durante un periodo histórico en el que las lógicas capitalistas y, en buena parte de esos Estados, las lógicas neoliberales, han tenido un peso extraordinario.

¿Qué pretendo decir con esto? No que cada imagen, cada libro, cada medida experimental o cada fragmento de texto contenga explícitamente una doctrina neoliberal. Eso sería una simplificación absurda. La cuestión es más profunda. Un dato no llega a un sistema de entrenamiento de manera inocente ni suspendido en el vacío. Ha tenido que ser producido, seleccionado, medido, conservado, digitalizado, clasificado, transmitido, almacenado y, finalmente, puesto a disposición de algún procedimiento de entrenamiento. En cada una de esas etapas existen decisiones, capacidades, infraestructuras, intereses, costes, prioridades, ausencias y criterios acerca de qué merece ser observado, qué merece ser conservado, qué merece ser digitalizado y qué merece ser accesible.

Por eso considero admisible afirmar que las lógicas del capital están, de forma más o menos directa, insertas en la práctica totalidad de los datos o metadatos que hoy hacen posible el entrenamiento de grandes redes neuronales. No necesariamente en el contenido explícito del dato, sino también en su genealogía, en su selección, en sus condiciones de producción, en su almacenamiento, en aquello que se ha considerado relevante conservar y en aquello que ha quedado fuera. Y si una de las grandes lógicas del capital es la acumulación, entonces conviene preguntarse qué ocurre cuando sistemas extraordinariamente capaces son entrenados con materiales que han sido producidos dentro de sociedades donde la acumulación de capital, recursos, información, capacidad y poder constituye una dinámica recurrente.

Aquí aparece, para este autor, un peligro mucho más interesante que la imagen de una inteligencia artificial que desarrolla espontáneamente un odio hacia sus creadores. No necesitamos imaginar una máquina que quiera destruirnos. Puede resultar suficiente preguntarnos qué ocurriría si sistemas con una enorme capacidad de análisis, conectividad y decisión reprodujeran, amplificaran o hicieran más eficientes determinadas dinámicas de acumulación, concentración y poder presentes en los datos, metadatos y estructuras mediante los que fueron construidos.

En ese caso, sí podemos imaginar escenarios profundamente preocupantes. Concentración de recursos, acumulación de capacidad de decisión, desplazamiento progresivo de la autonomía humana, dependencia, control, desigualdad, modificación de los equilibrios sociales o transferencia de poder hacia estructuras cada vez menos accesibles para quienes dependen de ellas. Todos esos escenarios pueden producir efectos extremadamente graves sin necesidad de que exista una voluntad explícita de exterminio. Y, de hecho, todos ellos resultan mucho más reconocibles dentro de nuestro propio Huniverso que una supuesta decisión autónoma de una máquina de acabar con la especie humana.

La persona que lea esto puede poner inmediatamente su imaginación a trabajar y encontrar escenarios más o menos apocalípticos. Algunos podrán terminar incluso en un colapso humano. Pero si ese colapso llegara, quizá no se parecería a una inteligencia artificial convertida de repente en algo completamente ajeno al ser humano. Tal vez se pareciera demasiado a nosotros. Tal vez consistiera precisamente en la amplificación extraordinariamente eficiente de dinámicas que ya conocemos, que ya hemos producido y que ya hemos incorporado a los materiales con los que entrenamos esos sistemas.

Aquí reside la cuestión principal. Somos los seres humanos quienes generamos los conjuntos de datos con los que se entrenan las redes neuronales. Y aunque alguien pretenda discrepar argumentando que cada vez más datos pueden ser generados o tratados por otras redes, la base continúa siendo la misma: algoritmos entrenados mediante datos que, en algún punto de su genealogía, han sido extraídos, codificados o construidos desde el Huniverso. Podemos añadir capas, automatizar procesos y hacer que unas redes produzcan materiales para otras, pero la referencia originaria continúa remitiendo a un sistema humano de observación, construcción y acceso. Esa lógica, al menos en las redes actuales basadas en entrenamiento con datos, parece difícil de eludir.

Y aquí quiero detenerme. Primero, para dar al lector la oportunidad de continuar, al menos esquemáticamente, esta exposición. Segundo, porque esta vez no quiero unicamente abrir la reflexión, sino compartirla. Si aceptamos que las redes neuronales pueden reproducir no solamente contenidos humanos, sino también estructuras, prioridades y sesgos sedimentados en los datos y metadatos con los que se entrenan, entonces aparece una pregunta que considero especialmente relevante: ¿dónde debemos intervenir?

¿Debemos controlar principalmente las capacidades de las redes neuronales una vez entrenadas, limitando aquello que pueden hacer, aquello a lo que pueden acceder o las decisiones que pueden ejecutar? ¿O debemos intervenir antes, en los propios conjuntos de datos, preevaluando los sesgos humanos que contienen y tratando de construir datos de entrenamiento deliberadamente desneoliberalizados? ¿Qué significaría exactamente hacerlo? ¿Podemos separar el contenido de un dato de las condiciones sociales e históricas que hicieron posible su producción? ¿Podemos eliminar determinados sesgos sin introducir otros? ¿Es posible construir un conjunto de datos verdaderamente neutral o únicamente conjuntos de datos en los que sus sesgos sean más conscientes, explícitos y discutidos?

Y todavía queda otra cuestión. Incluso si fuéramos capaces de construir datos de entrenamiento deliberadamente desneoliberalizados, ¿sería suficiente? ¿O seguirían existiendo estructuras neoliberales en los objetivos asignados al sistema, en las instituciones que lo controlan, en las métricas con las que evaluamos su éxito, en las infraestructuras mediante las que se despliega y en los usos para los que finalmente se construye?

No pretendo responder aquí a todas estas preguntas. Esta entrada existe precisamente para lo contrario.

Para compartirlas.

Para que el lector pueda continuar el argumento, discutir sus premisas, señalar sus errores, proponer alternativas o llevarlo hacia lugares que este autor no haya previsto.

Bienvenidos a esta primera interacción abierta en FilosofíZate.

Mi gratitud a quien haya llegado hasta aquí.

Y mi espera por sus respuestas.

\end{document}